% chapters/requirements/rf_grading.tex
% RF-11: Calificaciones
% ============================================================================

\item \textbf{Calificaciones.} \label{rf:grading}
Las calificaciones representan la valoración numérica de cada problema de una instancia de
examen. Son valores decimales que pueden ser negativos a nivel de problema, aunque la
calificación total del examen tiene como mínimo~0. La nota total se calcula siempre como la
suma de las calificaciones de los problemas. El sistema soporta la doble corrección
mediante concurrencia optimista: si dos correctores intentan calificar el mismo problema,
se compara la nota anterior conocida por el corrector con la almacenada y solo se aplica el
cambio si coinciden.

\begin{functional}

    % RF-11.1 ---------------------------------------------------------------
    \item \textbf{Asignación manual de nota a problema.} \label{rf:grading:manual}
    El sistema debe permitir a un corrector asignar o modificar la calificación numérica de
    un problema.

    \textbf{Entrada:} petición PUT al recurso de calificación del problema en la instancia,
    con la nota nueva y la nota anterior conocida por el corrector.

    \textbf{Proceso:} el sistema implementa concurrencia optimista: compara la nota anterior
    enviada por el corrector con la nota actualmente almacenada. Si coinciden, aplica la
    nueva nota. Si no coinciden (otro corrector la modificó entre tanto), rechaza la
    operación indicando que los datos están desactualizados y devuelve la calificación
    actual. Las calificaciones pueden ser decimales y negativas. Se registra la modificación
    en el \textit{log} de auditoría con la nota anterior, la nueva, el autor y la fecha.

    \textbf{Salida:} respuesta con la calificación actualizada, o error \acs{http}~409 si
    hay conflicto de concurrencia, incluyendo la calificación actual.

    % RF-11.2 ---------------------------------------------------------------
    \item \textbf{Asignación de nota por rúbrica.} \label{rf:grading:rubric}
    El sistema debe permitir a un corrector calificar un problema marcando o desmarcando los
    \textit{checkmarks} de su rúbrica.

    \textbf{Entrada:} petición PUT al recurso de rúbrica del problema en la instancia, con
    la lista de criterios marcados (identificadores de los \textit{checkmarks} activos).

    \textbf{Proceso:} el sistema calcula la calificación del problema como la suma de las
    puntuaciones de los criterios marcados (incluyendo criterios negativos como
    penalizaciones). Se aplica concurrencia optimista sobre el estado de la rúbrica. La
    calificación resultante se asigna al problema automáticamente.

    \textbf{Salida:} respuesta con el estado de la rúbrica (criterios marcados), la
    calificación resultante y, si procede, error de conflicto.

    % RF-11.3 ---------------------------------------------------------------
    \item \textbf{Cálculo automático de nota total.} \label{rf:grading:total}
    El sistema debe calcular la nota total de una instancia como la suma de las
    calificaciones de sus problemas, con un mínimo de~0.

    \textbf{Entrada:} evento interno disparado tras cualquier modificación de una
    calificación de problema.

    \textbf{Proceso:} el sistema recalcula la nota total de la instancia sumando las
    calificaciones de todos los problemas del modelo. Si la suma es negativa, la nota total
    se establece en~0. Los problemas sin calificar se consideran como~0 a efectos del
    cálculo, pero la instancia no se marca como completamente corregida hasta que todos los
    problemas tengan calificación explícita.

    \textbf{Salida:} nota total actualizada en la instancia.

    % RF-11.4 ---------------------------------------------------------------
    \item \textbf{Consulta de calificaciones desglosadas.} \label{rf:grading:detail}
    El sistema debe permitir consultar las calificaciones desglosadas de una instancia.

    \textbf{Entrada:} petición GET al recurso de calificaciones de la instancia.

    \textbf{Proceso:} el sistema recupera la calificación de cada problema, el estado de la
    rúbrica (si aplica), la nota total y el estado de corrección (completa o parcial). El
    nivel de detalle depende del rol del usuario y los permisos configurados en el examen:
    el corrector ve todo; el estudiante ve solo los datos permitidos (nota desglosada,
    rúbrica, según \cref{rf:exams:student_perms}).

    \textbf{Salida:} respuesta con las calificaciones desglosadas y la nota total.

    % RF-11.5 ---------------------------------------------------------------
    \item \textbf{Concurrencia optimista en calificaciones.}
    \label{rf:grading:concurrency}
    El sistema debe resolver conflictos de escritura simultánea en calificaciones mediante
    concurrencia optimista.

    \textbf{Entrada:} cualquier petición de modificación de calificación que incluya la nota
    anterior conocida por el corrector.

    \textbf{Proceso:} antes de aplicar la nueva calificación, el sistema compara la nota
    anterior enviada con la almacenada en base de datos. Si coinciden, la operación se
    ejecuta. Si difieren, la operación se rechaza con un código de conflicto, devolviendo
    la calificación actual para que el corrector pueda sincronizarse. Este mecanismo evita
    la pérdida silenciosa de calificaciones cuando varios correctores trabajan sobre el
    mismo problema.

    \textbf{Salida:} confirmación de la operación o error \acs{http}~409 con la
    calificación actual.

    % RF-11.6 ---------------------------------------------------------------
    \item \textbf{Consulta de nota por alumno.} \label{rf:grading:student}
    El sistema debe permitir al alumno consultar su nota una vez publicada.

    \textbf{Entrada:} petición GET al recurso de calificación de la instancia propia.

    \textbf{Proceso:} el sistema verifica que la instancia está en estado
    \texttt{PUBLISHED} o posterior (excluyendo \texttt{ARCHIVED}). Si el examen tiene
    activado el permiso de ver nota desglosada, se devuelve la calificación de cada
    problema y la nota total. Si solo está activado el permiso de ver la nota total, se
    omite el desglose. Si la instancia no está publicada, se deniega el acceso.

    \textbf{Salida:} respuesta con la nota (desglosada o total según permisos), o error
    \acs{http}~403 si la instancia no está publicada.

\end{functional}
